% ============================================================
\section{Remarks and open problems}\label{sec:remarks}
% ============================================================

\subsection{Consistency with the case $A=0$}
For $A=0$ we have $\sig_0=-2n_{0,2}$ and \eqref{eq:kappa-range} reads $0<\kappa<n_{0,2}$. The strip has normal width $2t\kappa$, i.e.\ vertical width $2t\kappa/n_{0,2}<2t$ (the curves are graphs).
This is consistent with the maximal vertical opening speed $2$ of the mixing zone for equal viscosities in our normalisation \cite{Ott99,Sze12,CCF21,FS18}.

\subsection{Optimal opening rates (formal)}
Single-strip subsolutions do not exhaust the relaxation. Consider instead a mixing zone of width $O(t)$ around $Z$ in which the density has a continuous profile $\bar\rho=R(\xi)$, $\xi=(x-Z)\cdot n/t$,
$R$ increasing from $-1$ at $\xi_-$ to $1$ at $\xi_+$, normalised so that $Z$ is the centre of mass, $\int_{\xi_-}^{\xi_+}\xi R'(\xi)\dd\xi=0$. Repeating the computations of Section~\ref{sec:proof} at leading
order (the conservation law now determines the normal component of $D$ by an integration across the zone instead of the boundary fluxes \eqref{eq:Fplus}--\eqref{eq:Fminus}), one finds, with
$\mathcal P(\xi):=-\int_{\xi_-}^{\xi}\eta R'(\eta)\dd\eta\ge0$ and $D\cdot\tau=0$,
\[
D\cdot n=-\sig-\frac{2(1-AR)\,\mathcal P}{1-R^2},\qquad a\cdot n=-\frac{\sig}{2}-\frac{A^2\mathcal P}{1+AR},
\]
and $D\cdot(D-2a)<0$ if and only if $(1-AR)\mathcal P<-\frac{\sig}2(1-R^2)$ on the zone. The extremal profile $\mathcal P=-\frac\sig2\frac{1-R^2}{1-AR}$ corresponds to $\xi=\frac{\sig}{2}F_1'(R)$ with
$F_1$ as in Example~\ref{ex:flat}, whose fronts are at
\[
\xi_+=-\frac{\sig}{1-A},\qquad \xi_-=\frac{\sig}{1+A}.
\]
Thus, formally, the upper and lower boundaries of the mixing zone may move away from the central interface with normal speeds up to $-\sig/(1-A)$ and $-\sig/(1+A)$, and the width may grow at any
rate below $-2\sig/(1-A^2)$. For the flat interface ($\sig=-2$) this is Example~\ref{ex:flat} with $\theta\to1$. By analogy with \cite[Section~5]{FS18}, we expect that piecewise constant subsolutions with
finitely many intermediate densities approximate these rates, and that the proof of Theorem~\ref{thm:main} extends to them with the same use of $(I-2A\Kop(z_0))^{-1}$.

\subsection{The role of the Saffman--Taylor function}
The proof shows that, at the level of subsolutions, the Saffman--Taylor function plays three roles at once: it is the density of the source part of the Darcy field (so that the field of the
subsolution is linear in $\sig$), it bounds the admissible opening rate, and it is the unknown of the only operator that has to be inverted. The gravitational and viscous contributions to the instability
enter only through $\sig=-2(n_2+AU)$, which is why the mixing zone of Example~\ref{ex:flat} is asymmetric: the relaxed flux is weighted by the mobility $1/(1-A\rho)$, in the spirit of the Koval model for miscible
viscous fingering \cite{Kov63,MO05}.

\subsection{Open problems}
\begin{enumerate}[label=(\arabic*),leftmargin=*]
\item \emph{Partially unstable data.} If $\sig_0$ changes sign, mixing should only occur where $\sig_0<0$, as in \cite{CFM22} for $A=0$. This requires $\kappa$ to vanish at the boundary of the unstable region,
which is excluded by the lower bound $\min\kappa>0$ used in Lemma~\ref{lem:stream}.
\item \emph{Longer times.} The time expansion only gives a short existence time. Continuing the construction requires an evolution equation for the central interface, as in \cite{CCF21,ACF21}; when $A\neq0$
it should involve the averaged version of $(I-2A\Kop)^{-1}$ along the mixing zone.
\item \emph{Regularity.} We assumed $z_0\in C^{6,\gamma}$, which is far from the $C^{3,\alpha}$ of \cite{FS18} for $A=0$. The losses come from the crude bookkeeping in Lemma~\ref{lem:timereg}.
\item \emph{Selection.} For $A=0$, Otto's relaxation selects the maximal mixing rate \cite{Ott99,Sze12}; it would be interesting to know whether a similar selection principle singles out the rates of the previous
subsection when $A\neq0$, see \cite{Men22} for the flat case.
\end{enumerate}
